%% 12_cmcontrol.tex -- Chapter 12: Common-Mode Control and Electrode Monitoring
%% Status: OUTLINE (build stage 3). Replace the outline with the full chapter when it is written.
\chapter{Common-Mode Control and Electrode Monitoring}\label{ch:cmcontrol}
\chapterlead{The third electrode is not a reference: it is the output of a feedback loop that pins the body to the amplifier's common mode.}

\begin{specnote}[Chapter status: outline]
This chapter is an outline. It will be written in build stage 3 (Part II with Labs 3--5). The topics, the design question and the labs below are fixed; the derivations, worked examples and figures are still to come.
\end{specnote}

\begin{objectives}
\begin{itemize}
  \item design a right-leg drive (driven-right-leg) loop and analyse its stability;
  \item relate RLD loop gain to the residual common-mode and differential interference;
  \item implement common-mode feedback for the multichannel front-end;
  \item detect lead-off with DC and AC methods;
  \item protect the inputs against ESD and defibrillation pulses in \sgx{}.
\end{itemize}
\end{objectives}

\begin{insight}[Design question]
How much right-leg-drive gain and which protection network let the \OB{} channel pass the ECG CMRR test with margin?
\end{insight}

\section{The right-leg drive loop}\label{sec:ch12-1}
Sensing the average input common mode and driving it back into the body; the loop gain reduces the body common-mode voltage of \cref{ch:interference}; stability with the electrode and body capacitances~\cite{winter1983}.

\section{Common-mode feedback}\label{sec:ch12-2}
Setting the output common mode of the fully differential stages across eight channels.

\section{Lead-off detection}\label{sec:ch12-3}
DC methods (small pull-up/pull-down currents) and AC methods (an excitation above the signal band); leakage-current limits from \cref{ch:interference}; the full impedance measurement follows in \cref{ch:bioz}.

\section{Input protection and ESD}\label{sec:ch12-4}
Series resistors, clamps and their noise and leakage cost; defibrillation-proof inputs; the ESD devices available in \sgx{}.

\begin{tryit}[Labs for this chapter]
Lab~2 (\cref{sec:lab2}) already adds a behavioural right-leg drive to the body model of \cref{ch:interference}.
\end{tryit}

\paragraph{Read with this chapter} Winter and Webster~\cite{winter1983}, Huhta and Webster~\cite{huhta1973}, Webster~\cite{webster}.
